\section{Diagnostiche e risultati per funzione}
\label{app:extended-results}

\subsection{Audit sincrono e variabilità interna}
La Figura~\ref{fig:result-sync-equivalence} rappresenta le 18 righe del riepilogo di equivalenza. Le tre righe per combinazione backend--worker corrispondono ai gruppi da 60, 10 e 10 confronti; la Tabella~\ref{tab:result-equivalence} ne riporta la somma. La Figura~\ref{fig:app-internal-forest} conserva il dettaglio degli effetti di qualità per funzione e dimensione rispetto ai riferimenti seriali appropriati. Le differenze aggregate del corpo principale vanno lette insieme a questa dispersione.

\begin{figure}[H]
\centering
\includegraphics[width=0.68\textwidth]{figures/results/fig_sync_equivalence.pdf}
\caption{Uguaglianza esatta dei campi registrati nei 18 gruppi dell\textquotesingle audit. Tutti i gruppi raggiungono il 100\%; le righe con etichetta ripetuta sono gruppi distinti nel riepilogo disponibile.}
\label{fig:result-sync-equivalence}
\end{figure}

\begin{figure}[H]
\centering
\includegraphics[width=\linewidth]{figures/results/fig_internal_function_forest.pdf}
\caption{Effetti appaiati per funzione e dimensione. Gli intervalli bootstrap al 95\% ricampionano i semi entro cella e sono condizionati alle funzioni analizzate.}
\label{fig:app-internal-forest}
\end{figure}

\clearpage
\subsection{Sensibilità al passo temporale}
\label{app:dt-sensitivity-results}
Lo screening comprende 900 run: sei funzioni, due dimensioni, quindici valori di $\Delta t$ e cinque semi, con $N=64$, $K=500$ e 32\,064 FE/run. Le differenze sono appaiate rispetto a $\Delta t=30$. La Tabella~\ref{tab:app-dt-sensitivity} aggrega le sei mediane per funzione a ogni dimensione; la Figura~\ref{fig:app-dt-sensitivity} conserva il dettaglio delle singole funzioni. Il riferimento produce differenze nulle per costruzione.

Le risposte al passo sono eterogenee tra funzioni e dimensioni. Diversi valori mostrano mediane aggregate negative rispetto al riferimento, con entità variabile; lo screening descrive questa sensibilità senza identificare un unico passo valido per tutti i problemi. I quartili in tabella descrivono la variabilità tra le sei mediane per funzione.

\begin{table}[H]
\centering
\caption{Mediana [Q1, Q3] tra le sei mediane per funzione delle differenze appaiate di log-errore rispetto a $\Delta t=30$. Le 60 run per passo comprendono entrambe le dimensioni.}
\label{tab:app-dt-sensitivity}
\small
\begin{tabular}{@{}rllr@{}}
\toprule
$\Delta t$ & $D=30$ & $D=100$ & Run \\
\midrule
0.1 & -0.971 [-1.487, -0.492] & -0.456 [-0.688, -0.270] & 60 \\
0.2 & -1.122 [-1.712, -0.582] & -0.510 [-0.817, -0.303] & 60 \\
0.35 & -1.053 [-1.667, -0.457] & -0.484 [-0.761, -0.264] & 60 \\
0.5 & -1.074 [-1.793, -0.497] & -0.493 [-0.751, -0.249] & 60 \\
0.7 & -0.895 [-1.450, -0.462] & -0.482 [-0.729, -0.232] & 60 \\
1 & -0.777 [-1.167, -0.402] & -0.434 [-0.724, -0.228] & 60 \\
1.5 & -0.423 [-0.608, -0.272] & -0.293 [-0.493, -0.159] & 60 \\
2.2 & -0.491 [-0.733, -0.257] & -0.288 [-0.487, -0.123] & 60 \\
3.6 & -0.411 [-0.571, -0.221] & -0.304 [-0.438, -0.124] & 60 \\
5.5 & -0.251 [-0.396, -0.095] & -0.154 [-0.291, -0.065] & 60 \\
8 & -0.180 [-0.224, -0.074] & -0.119 [-0.187, -0.072] & 60 \\
12 & -0.128 [-0.159, -0.079] & -0.053 [-0.103, -0.026] & 60 \\
20 & -0.059 [-0.087, -0.036] & -0.031 [-0.057, -0.020] & 60 \\
30 & 0.000 [0.000, 0.000] & 0.000 [0.000, 0.000] & 60 \\
50 & 0.023 [0.020, 0.026] & 0.006 [-0.001, 0.021] & 60 \\
\bottomrule
\end{tabular}
\end{table}

\begin{figure}[H]
\centering
\includegraphics[width=\linewidth]{figures/results/fig_dt_sensitivity_summary.pdf}
\caption{Sensibilità del log-errore finale al passo, per funzione e dimensione. Valori negativi indicano errori inferiori al riferimento $\Delta t=30$.}
\label{fig:app-dt-sensitivity}
\end{figure}

\FloatBarrier
\subsection{Migrazioni per funzione}
La Tabella~\ref{tab:app-migration-functions} riporta i conteggi per funzione. Ogni cella usa cinque run da 36 tentativi, per un totale di 180 tentativi; le percentuali aggregate del corpo principale sommano questi conteggi su tutte le funzioni. La Figura~\ref{fig:app-migration-diagnostics} mostra la distribuzione dei tassi di accettazione nelle configurazioni con migrazione.

\begin{table}[H]
\centering
\caption{Migrazioni accettate/tentate, per funzione. G--R: vicinato global e migrazione ring; G--C: global e migrazione completa; R--R: ring locale e migrazione ring.}
\label{tab:app-migration-functions}
\small
\begin{tabular}{@{}lrrrrrr@{}}
\toprule
& \multicolumn{2}{c}{G--R} & \multicolumn{2}{c}{G--C} & \multicolumn{2}{c}{R--R} \\ Funzione & $D=30$ & $D=100$ & $D=30$ & $D=100$ & $D=30$ & $D=100$ \\
\midrule
Sphere & 99/180 & 107/180 & 158/180 & 153/180 & 158/180 & 160/180 \\
Rastrigin & 157/180 & 161/180 & 168/180 & 169/180 & 178/180 & 177/180 \\
Rosenbrock & 100/180 & 103/180 & 155/180 & 156/180 & 149/180 & 168/180 \\
Ackley & 136/180 & 135/180 & 177/180 & 164/180 & 179/180 & 169/180 \\
Griewank & 99/180 & 107/180 & 158/180 & 153/180 & 158/180 & 160/180 \\
Salomon & 172/180 & 164/180 & 179/180 & 179/180 & 179/180 & 178/180 \\
\bottomrule
\end{tabular}
\end{table}

\begin{figure}[H]
\centering
\includegraphics[width=\linewidth]{figures/results/fig_migration_diagnostics.pdf}
\caption{Diagnostica della migrazione. I tentativi sono determinati da epoche e numero di isole; le sostituzioni accettate provengono dai conteggi registrati.}
\label{fig:app-migration-diagnostics}
\end{figure}

\FloatBarrier
\subsection{BBOB: ranghi mediani per funzione}
I ranghi sono calcolati entro la stessa istanza e poi aggregati sulle cinque istanze. Le Tabelle~\ref{tab:app-bbob-d5} e~\ref{tab:app-bbob-d20} riportano le mediane disponibili; la Figura~\ref{fig:app-external-heatmaps} ne offre una lettura congiunta. Si tratta di ranghi finali, distinti dalle frazioni di target raggiunti della Tabella~\ref{tab:result-bbob}.

\begin{table}[H]
\centering
\caption{BBOB, $D=5$: mediana dei ranghi entro istanza su cinque istanze per funzione.}
\label{tab:app-bbob-d5}
\small
\begin{tabular}{@{}lrrrrrrr@{}}
\toprule
Metodo & F1 & F5 & F8 & F13 & F17 & F21 & F24 \\
\midrule
Classic gbest & 11.0 & 4.5 & 11.0 & 10.0 & 10.0 & 8.0 & 10.0 \\
CLPSO & 10.0 & 11.0 & 10.0 & 11.0 & 11.0 & 4.0 & 10.0 \\
Constriction gbest & 4.0 & 4.5 & 2.0 & 6.0 & 3.0 & 10.0 & 3.0 \\
CPSO & 6.0 & 4.5 & 8.0 & 7.0 & 7.0 & 10.0 & 4.0 \\
CSO & 1.5 & 9.0 & 6.0 & 2.0 & 1.0 & 6.0 & 3.0 \\
DMS-PSO & 6.0 & 4.5 & 5.0 & 4.0 & 5.0 & 3.0 & 5.0 \\
FIPS & 9.0 & 4.5 & 9.0 & 8.0 & 9.0 & 5.0 & 9.0 \\
Inertia gbest & 5.0 & 4.5 & 3.0 & 7.0 & 4.0 & 8.0 & 7.0 \\
Inertia ring & 8.0 & 4.5 & 4.0 & 5.0 & 7.0 & 3.0 & 3.0 \\
SL-PSO & 1.5 & 4.5 & 5.0 & 5.0 & 6.0 & 8.0 & 5.0 \\
SPSO-2011 & 3.0 & 10.0 & 2.0 & 1.0 & 3.0 & 2.0 & 6.0 \\
\bottomrule
\end{tabular}
\end{table}

\begin{table}[H]
\centering
\caption{BBOB, $D=20$: mediana dei ranghi entro istanza su cinque istanze per funzione.}
\label{tab:app-bbob-d20}
\small
\begin{tabular}{@{}lrrrrrrr@{}}
\toprule
Metodo & F1 & F5 & F8 & F13 & F17 & F21 & F24 \\
\midrule
Classic gbest & 11.0 & 8.0 & 11.0 & 11.0 & 9.0 & 11.0 & 10.0 \\
CLPSO & 9.0 & 9.0 & 10.0 & 10.0 & 10.0 & 9.0 & 11.0 \\
Constriction gbest & 2.5 & 2.5 & 2.0 & 4.0 & 8.0 & 5.0 & 4.0 \\
CPSO & 10.0 & 11.0 & 9.0 & 9.0 & 11.0 & 8.0 & 8.0 \\
CSO & 2.5 & 6.0 & 3.0 & 2.0 & 1.0 & 4.0 & 1.0 \\
DMS-PSO & 6.0 & 2.5 & 4.0 & 5.0 & 3.0 & 3.0 & 5.0 \\
FIPS & 8.0 & 7.0 & 7.0 & 8.0 & 6.0 & 6.0 & 9.0 \\
Inertia gbest & 5.0 & 2.5 & 3.0 & 4.0 & 5.0 & 6.0 & 6.0 \\
Inertia ring & 7.0 & 5.0 & 5.0 & 7.0 & 6.0 & 1.0 & 6.0 \\
SL-PSO & 2.5 & 2.5 & 6.0 & 3.0 & 7.0 & 7.5 & 2.0 \\
SPSO-2011 & 2.5 & 10.0 & 3.0 & 2.0 & 2.0 & 7.0 & 4.0 \\
\bottomrule
\end{tabular}
\end{table}

\begin{figure}[H]
\centering
\includegraphics[width=\linewidth]{figures/results/fig_external_function_heatmaps.pdf}
\caption{Ranghi mediani BBOB per funzione e dimensione. Il rango 1 corrisponde al valore obiettivo più basso; metodi in ordine alfabetico, senza evidenziazioni selettive.}
\label{fig:app-external-heatmaps}
\end{figure}

\FloatBarrier
\subsection{CEC: confronto per funzione}
La Figura~\ref{fig:app-cec-breakdown} riporta il rango medio sui tre semi a 120 mila FE per ciascuna funzione. Il dettaglio evidenzia come l'ordinamento cambi tra funzioni, anche all'interno dello stesso gruppo dimensionale. L'Appendice~\ref{app:cec-values} riporta i valori obiettivo corrispondenti ai tre checkpoint, mantenendo le scale delle funzioni separate.

\begin{figure}[H]
\centering
\includegraphics[width=0.90\textwidth]{figures/results/fig_cec_function_breakdown.pdf}
\caption{CEC: rango medio entro funzione--seme a 120 mila FE. Ordine alfabetico dei metodi; F13 e F14 hanno 905 variabili, le altre quattro funzioni 1000.}
\label{fig:app-cec-breakdown}
\end{figure}
